\section{Replacement regions and the recursive repair family}
\label{sec:area_law_repair_regions}

The replacement construction of \cite[Section~11]{OpenAI2026AreaLaw}
uses the eight triangles of a square divided at its side midpoints.
We retain the removal half-side $d(S)$ and square $Q(v,S)$ of
Definition~\ref{def:area_law_repair_square}. Distances in the plane are
maximum-norm distances, and colors are represented by $0$ and $1$.

% Source: 10-geometry.tex, lines 393--417; geometry:repair-induction,
% lines 471--514; open actual regions and closed birth regions, lines 558--561.
% Revision: adc7f1241b42e322a6451854ab7e4b4c146bf78a.
% Completion marks below refer to the definitions, iteration and local
% disjointness lemma. The full repair induction and exceptional remainder
% are separate assertions.

\begin{definition}[The normalized fan and its colored runs]
    \label{def:area_law_repair_fan}
    \lean{TNLean.PEPS.AreaLaw.Geometry.repairFanSplit,
        TNLean.PEPS.AreaLaw.Geometry.RepairFanSlot,
        TNLean.PEPS.AreaLaw.Geometry.repairExteriorColor,
        TNLean.PEPS.AreaLaw.Geometry.RepairFanRun,
        TNLean.PEPS.AreaLaw.Geometry.repairFanRunExteriorColor,
        TNLean.PEPS.AreaLaw.Geometry.repairReplacementColor}
    \leanok
    \uses{def:area_law_fan_runs,thm:area_law_fan_run_color}
    Dividing each side of $[-1,1]^2$ at its midpoint gives $8$ cyclically ordered perimeter
    segments $[a_s,b_s]$ indexed by $J=\{0,1,2,3\}\times\{0,1\}$, with closed fan triangles
    $T_s=\operatorname{conv}\{0,a_s,b_s\}$. Given region identifiers $I$, colors $\xi:I\to\{0,1\}$,
    and an assignment $\sigma:J\to I$, set the exterior color $\lambda_\sigma(s)=\xi(\sigma(s))$
    and let $\mathcal R(\xi,\sigma)$ be the connected components (runs, Definition~\ref{def:area_law_fan_runs})
    of the graph joining consecutive $s\in J$ with equal $\lambda_\sigma(s)$. For each
    $H\in\mathcal R(\xi,\sigma)$ and any $s\in H$ (by Lemma~\ref{thm:area_law_fan_run_color}), define
    \begin{align}
        T_H & =\bigcup_{s\in H}T_s,
            & \lambda_\sigma(H)                & =\lambda_\sigma(s),
            & \lambda_\sigma^{\mathrm{new}}(H) & =1-\lambda_\sigma(H).
        \label{eq:area_law_repair_run}
    \end{align}
\end{definition}

% Source: 10-geometry.tex, replacement instructions, lines 404--417;
% open actual regions and closed birth regions, lines 558--561.
\begin{definition}[Birth and current region of a replacement]
    \label{def:area_law_repair_birth}
    \lean{TNLean.PEPS.AreaLaw.Geometry.repairAffineMap,
        TNLean.PEPS.AreaLaw.Geometry.repairReplacementBirthRegion,
        TNLean.PEPS.AreaLaw.Geometry.repairReplacementOpenRegion}
    \leanok
    \uses{def:area_law_repair_fan,def:area_law_repair_square}
    For a center $v\in\mathbb R^2$, scale $S\in\mathbb R$, and run $H\in\mathcal R(\xi,\sigma)$, define
    \begin{align}
        F_{v,S}(x) & =v+d(S)x,
                   & B(v,S,H)  & =F_{v,S}(T_H),
                   & U(v,S,H)  & =\operatorname{int}B(v,S,H).
        \label{eq:area_law_repair_birth}
    \end{align}
\end{definition}

% Source: 10-geometry.tex, lines 404--417 and 486--514;
% closed births and open current regions, lines 558--561.
\begin{lemma}[Disjoint replacement regions at one parent]
    \label{thm:area_law_repair_replacement_disjoint}
    \lean{TNLean.PEPS.AreaLaw.Geometry.repairReplacementOpenRegion_pairwise_disjoint}
    \leanok
    \uses{def:area_law_repair_birth,thm:area_law_fan_run_disjoint_interiors}
    For any $v\in\mathbb R^2$, $S\in\mathbb R$, and exterior assignment $\sigma$, distinct runs
    $H\ne K$ satisfy $U(v,S,H)\cap U(v,S,K)=\varnothing$.
\end{lemma}
\begin{proof}\leanok
    Since $d(S)>0$, $F_{v,S}(x)=v+d(S)x$ is a homeomorphism preserving interiors and disjointness;
    apply Lemma~\ref{thm:area_law_fan_run_disjoint_interiors} to $\operatorname{int}T_H$ and $\operatorname{int}T_K$.
\end{proof}

% Source: 10-geometry.tex, replacement sectors and their triangle unions,
% lines 404--417; translated supporting lines and lattice sites, lines 545--561.
% Revision: adc7f1241b42e322a6451854ab7e4b4c146bf78a.
\begin{lemma}[Lattice points in a replacement square]
    \label{thm:area_law_repair_replacement_lattice_cover}
    \lean{TNLean.PEPS.AreaLaw.Geometry.exists_mem_repairReplacementOpenRegion_of_mem_repairSquare}
    \leanok
    \uses{def:area_law_dyadic_origin,def:area_law_affine_mesh,
        def:area_law_repair_birth}
    For $o=(\sqrt2,\sqrt3)$, $v\in\mathcal M(o,1)=o+\mathbb Z^2$, $S\ge1$, and any $I$,
    $\xi:I\to\{0,1\}$, and $\sigma:J\to I$,
    \begin{align}
        Q(v,S)\cap\mathbb Z^2
         & \subseteq\bigcup_{H\in\mathcal R(\xi,\sigma)}U(v,S,H).
        \label{eq:area_law_repair_replacement_lattice_cover}
    \end{align}
\end{lemma}

\begin{proof}
    \leanok
    \uses{def:area_law_repair_square,thm:area_law_repair_halfside_integral,
        thm:area_law_dyadic_origin_lines,lem:area_law_repair_fan_cones,
        def:area_law_repair_birth}
    Let $v=o+m$ ($m\in\mathbb Z^2$) and $d=d(S)\in\mathbb Z_{>0}$
    (Lemma~\ref{thm:area_law_repair_halfside_integral}). For $z\in Q(v,S)\cap\mathbb Z^2$,
    $\partial Q(v,S)$ and all fan rays lie on integer translates of the four line families of
    Lemma~\ref{thm:area_law_dyadic_origin_lines}, which miss $\mathbb Z^2$, so $\|z-v\|_\infty<d$
    and $z-v$ avoids all fan rays. Lemma~\ref{lem:area_law_repair_fan_cones} then places
    $z\in\operatorname{int}(v+dT_s)\subseteq\operatorname{int}B(v,S,H)=U(v,S,H)$ for the run $H\ni s$.
\end{proof}

% Source: 10-geometry.tex, lines 404--407; geometry:repair-induction,
% working square and current exterior sectors, lines 486--504.
% Existence and uniqueness of a containing current region remain to be proved
% for the actual recursively constructed family.
\begin{definition}[Assignments from the current working wedges]
    \label{def:area_law_repair_current_assignment}
    \lean{TNLean.PEPS.AreaLaw.Geometry.repairCurrentFanTriangle,
        TNLean.PEPS.AreaLaw.Geometry.repairCurrentSectorAssignment,
        TNLean.PEPS.AreaLaw.Geometry.repairCurrentAssignments}
    \leanok
    \uses{def:area_law_repair_fan,def:area_law_repair_square}
    For a nonempty identifier set $I$ with current regions $(R_i)_{i\in I}$, $v\in\mathbb R^2$,
    $r\in\mathbb R$, and $s\in J$, set $W_s(v,r)=v+rT_s$ and define
    \begin{align}
        \sigma^R_{v,r}(s)
         & =\varepsilon\{i\in I:\operatorname{int}W_s(v,r)\subseteq R_i\}, &
        \sigma^R_p(s)
         & =\sigma^R_{v_p,\,6d(S_p)}(s),
        \label{eq:area_law_repair_assignment}
    \end{align}
    where $\varepsilon$ chooses an element of the set when nonempty and an arbitrary element of $I$ otherwise.
\end{definition}

% Source: 10-geometry.tex, lines 393--409; geometry:repair-induction,
% lines 471--514; closed births and current regions, lines 558--561.
\begin{definition}[Simultaneous replacement of the current family]
    \label{def:area_law_repair_successor_regions}
    \lean{TNLean.PEPS.AreaLaw.Geometry.RepairReplacementIndex,
        TNLean.PEPS.AreaLaw.Geometry.RepairRegionIndex,
        TNLean.PEPS.AreaLaw.Geometry.repairNextBirthRegion,
        TNLean.PEPS.AreaLaw.Geometry.repairNextCurrentRegion,
        TNLean.PEPS.AreaLaw.Geometry.repairNextRegionColor}
    \leanok
    \uses{def:area_law_repair_birth,def:area_law_simultaneous_remnants}
    At generation $j\ge0$, given $(B_i,R_i,\xi_i)_{i\in I}$, parent marks $K$ with $(v_p,S_p,\sigma_p)$,
    and stopping scale $S_*\in\mathbb R$, set
    \begin{align}
        K_+ & =\{p\in K:S_*\le d(S_p)\},
            & I_{\mathrm{new}}           & =\{j+1\}\times\coprod_{p\in K_+}\mathcal R(\xi,\sigma_p), \\
        I'  & =I\sqcup I_{\mathrm{new}},
            & Q_K                        & =\bigcup_{p\in K}Q(v_p,S_p),
        \label{eq:area_law_repair_indices}
    \end{align}
    and for $i\in I$ and $q=(j+1,p,H)\in I_{\mathrm{new}}$, define
    \begin{align}
        B'_i & =B_i,
             & R'_i           & =R_i\setminus Q_K,
             & \xi'_i         & =\xi_i,                   \\
        B'_q & =B(v_p,S_p,H),
             & R'_q           & =U(v_p,S_p,H),
             & \xi'_q         & =1-\lambda_{\sigma_p}(H).
        \label{eq:area_law_repair_successor}
    \end{align}
\end{definition}

% Source: OpenAI2026AreaLaw, Section 11, 10-geometry.tex,
% replacement instructions, lines 393--417; geometry:repair-induction,
% lines 478 and 486--514; open current regions, lines 558--561.
% Revision: adc7f1241b42e322a6451854ab7e4b4c146bf78a.
\begin{lemma}[Disjointness after simultaneous repair]
    \label{thm:area_law_repair_next_disjoint}
    \lean{TNLean.PEPS.AreaLaw.Geometry.repairNextCurrentRegion_pairwise_disjoint}
    \leanok
    \uses{def:area_law_repair_successor_regions}
    In Definition~\ref{def:area_law_repair_successor_regions}, if:
    (1) $R_i\cap R_k=\varnothing$ for distinct $i,k\in I$;
    (2) $Q(v_p,S_p)\cap Q(v_q,S_q)=\varnothing$ for distinct $p,q\in K$; and
    (3) $U(v_p,S_p,H)\subseteq Q(v_p,S_p)$ for all $p\in K_+$ and $H\in\mathcal R(\xi,\sigma_p)$,
    then $R'_x\cap R'_y=\varnothing$ for all distinct $x,y\in I'$.
\end{lemma}

\begin{proof}\leanok
    \uses{def:area_law_repair_successor_regions,
        thm:area_law_repair_replacement_disjoint}
    By (\ref{eq:area_law_repair_successor}): for distinct $i,k\in I$, $R'_i\cap R'_k\subseteq R_i\cap R_k=\varnothing$
    by (1); for $i\in I$ and $q=(j+1,p,H)\in I_{\mathrm{new}}$, $R'_q\subseteq Q(v_p,S_p)\subseteq Q_K$ by (3)
    while $R'_i=R_i\setminus Q_K$; for fresh identifiers with distinct parents $p\ne p'$,
    $R'_q\cap R'_{q'}\subseteq Q(v_p,S_p)\cap Q(v_{p'},S_{p'})=\varnothing$ by (2)--(3); and for
    distinct runs $H\ne H'$ at the same parent $p$, $U(v_p,S_p,H)\cap U(v_p,S_p,H')=\varnothing$
    by Lemma~\ref{thm:area_law_repair_replacement_disjoint}.
\end{proof}

% Source: 10-geometry.tex, terminal cuts and child marks, lines 398--409;
% geometry:repair-induction, lines 506--514.
\begin{definition}[Terminal interiors and next-generation marks]
    \label{def:area_law_repair_terminal_children}
    \lean{TNLean.PEPS.AreaLaw.Geometry.repairTerminalInteriorSet,
        TNLean.PEPS.AreaLaw.Geometry.RepairNextMarkIndex,
        TNLean.PEPS.AreaLaw.Geometry.repairNextMarkLocation,
        TNLean.PEPS.AreaLaw.Geometry.repairNextMarkScale}
    \leanok
    \uses{def:area_law_repair_successor_regions,def:area_law_repair_children}
    In the setting of Definition~\ref{def:area_law_repair_successor_regions}, define
    \begin{align}
        E_{\mathrm{term}}
                            & =\bigcup_{\substack{p\in K    \\d(S_p)<S_*}}\operatorname{int}Q(v_p,S_p), &
        P                   & =\{-1,0,1\}^2,                \\
        K'                  & =\{j+1\}\times K_+\times P, &
        v'_{(j+1,p,\omega)} & =v_p+d(S_p)\omega,          &
        S'_{(j+1,p,\omega)} & =d(S_p).
        \label{eq:area_law_repair_next_marks}
    \end{align}
\end{definition}

% Source: 10-geometry.tex, simultaneous removals, terminal interiors and
% replacements, lines 393--417; supporting-line exclusion, lines 545--561.
% Revision: adc7f1241b42e322a6451854ab7e4b4c146bf78a.
\begin{lemma}[Lattice coverage after one repair]
    \label{thm:area_law_repair_next_lattice_cover}
    \lean{TNLean.PEPS.AreaLaw.Geometry.exists_mem_repairNextCurrentRegion_or_terminalInterior}
    \leanok
    \uses{def:area_law_repair_successor_regions,
        def:area_law_repair_terminal_children,def:area_law_affine_mesh,
        def:area_law_dyadic_origin}
    In Definitions~\ref{def:area_law_repair_successor_regions} and
    \ref{def:area_law_repair_terminal_children} with $o=(\sqrt2,\sqrt3)$, if $v_p\in\mathcal M(o,1)$
    and $S_p\ge1$ for all $p\in K$, then
    \begin{align}
        \Bigl(\bigcup_{i\in I}R_i\Bigr)\cap\mathbb Z^2
         & \subseteq\Bigl(\bigcup_{x\in I'}R'_x\Bigr)\cup E_{\mathrm{term}}.
        \label{eq:area_law_repair_next_lattice_cover}
    \end{align}
\end{lemma}

\begin{proof}
    \leanok
    \uses{def:area_law_repair_successor_regions,
        def:area_law_repair_terminal_children,def:area_law_repair_birth,
        thm:area_law_repair_replacement_lattice_cover}
    For $z\in R_i\cap\mathbb Z^2$: if $z\notin Q_K$, then $z\in R_i\setminus Q_K=R'_i$;
    if $z\in Q(v_p,S_p)$ ($p\in K$), then (\ref{eq:area_law_repair_replacement_lattice_cover}) of
    Lemma~\ref{thm:area_law_repair_replacement_lattice_cover} places $z\in U(v_p,S_p,H)$ for some
    run $H$, which equals $R'_{(j+1,p,H)}$ when $d(S_p)\ge S_*$ and lies in
    $\operatorname{int}Q(v_p,S_p)\subseteq E_{\mathrm{term}}$ when $d(S_p)<S_*$.
\end{proof}

% Source: 10-geometry.tex, initial choices, lines 212--233 and 325--363;
% replacement instructions, lines 393--417; geometry:repair-induction,
% lines 471--514; current regions and remembered births, lines 558--561.
\begin{definition}[The recursive repair family]
    \label{def:area_law_repair_stage}
    \lean{TNLean.PEPS.AreaLaw.Geometry.RepairStage,
        TNLean.PEPS.AreaLaw.Geometry.initialRepairStage,
        TNLean.PEPS.AreaLaw.Geometry.repairNextStage,
        TNLean.PEPS.AreaLaw.Geometry.repairStage}
    \leanok
    \uses{def:area_law_repair_current_assignment,
        def:area_law_repair_successor_regions,
        def:area_law_repair_terminal_children,
        thm:area_law_uniform_separated_initial_marks,
        def:area_law_initial_open_regions,def:area_law_initial_region_colors}
    A stage is a tuple $\mathcal D=(I,*,B,R,\xi,K,v,S)$ with identifiers $I$, dummy $*\in I$,
    sets $B_i,R_i\subseteq\mathbb R^2$, labels $\xi_i\in\{0,1\}$, marks $K$, locations $v:K\to\mathbb R^2$,
    and scales $S:K\to\mathbb R$. Given finite $\Lambda$, $A\subseteq\Lambda$, $o\in\mathbb R^2$,
    $C_0\ge2$, $k_0\ge50{,}000{,}000$, and the jointly chosen sparse residues $(a_k,b_k)$, initial
    marks $M$, and scale function $s$ from Lemma~\ref{thm:area_law_uniform_separated_initial_marks}
    (inducing $\mathcal I$, $X_i^0$, and $\lambda(i)$ via
    Definitions~\ref{def:area_law_initial_open_regions} and \ref{def:area_law_initial_region_colors}),
    set $\mathcal D_{\mathrm{init}}$ by
    \begin{align}
        I_0       & =\mathcal I,
                  & *_0          & =*,
                  & K_0          & =M,                       \\
        B_i^{(0)} & =X_i^0,
                  & R_i^{(0)}    & =\operatorname{int}X_i^0,
                  & \xi_i^{(0)}  & =\lambda(i),              \\
        v_p^{(0)} & =p,
                  & S_p^{(0)}    & =s(p).
        \label{eq:area_law_repair_initial}
    \end{align}
    For $S_*\in\mathbb R$, let $\Phi_{j,S_*}(\mathcal D)$ apply (\ref{eq:area_law_repair_assignment}),
    (\ref{eq:area_law_repair_successor}), and (\ref{eq:area_law_repair_next_marks}), and define
    \begin{align}
        \mathcal D_0 & =\mathcal D_{\mathrm{init}},
                     & \mathcal D_{j+1}             & =\Phi_{j,S_*}(\mathcal D_j).
        \label{eq:area_law_repair_iteration}
    \end{align}
\end{definition}

% Source: 10-geometry.tex, geometry:belt-count and geometry:initial-stars,
% lines 212--233 and 325--363.
% Revision: adc7f1241b42e322a6451854ab7e4b4c146bf78a.
\begin{lemma}[Properties of the chosen initial repair marks]
    \label{thm:area_law_initial_repair_mark_facts}
    \lean{TNLean.PEPS.AreaLaw.Geometry.initialRepairStage_mark_spec}
    \leanok
    \uses{def:area_law_repair_stage,def:area_law_belt_cells,
        def:area_law_belt_cell_marks,def:area_law_fine_belt_scales}
    In Definition~\ref{def:area_law_repair_stage}, let $F_k$ be the selected belt-cell indices
    (Definition~\ref{def:area_law_belt_cells}) and $t_k=2^{\ell_k}$
    (Definition~\ref{def:area_law_fine_belt_scales}). Then:
    \begin{enumerate}
        \item $|F_k|\le64(2C_0+1)^2|\partial_\Lambda A|\,2^{-\delta_0 k/2}$ for all $k\ge0$, and
              $\{k\ge0:F_k\ne\varnothing\}$ is finite;
        \item for every $p\in K_0$, $S_p^{(0)}>0$, there exist $k\ge k_0$ and $z\in F_k$ with
              $v_p^{(0)}\in M_{\ell_k,z}$ (Definition~\ref{def:area_law_belt_cell_marks}) and
              $S_p^{(0)}=t_k$, and any $k\ge k_0$, $z\in F_k$ with $v_p^{(0)}\in M_{\ell_k,z}$
              satisfies $S_p^{(0)}\le t_k$;
        \item $p\mapsto v_p^{(0)}$ is injective, and distinct $p,q\in K_0$ satisfy
              $\operatorname{dist}_\infty(v_p^{(0)},v_q^{(0)})\ge\frac14\max(S_p^{(0)},S_q^{(0)})$.
    \end{enumerate}
\end{lemma}

\begin{proof}\leanok
    \uses{thm:area_law_uniform_separated_initial_marks,
        def:area_law_repair_stage}
    Immediate from Lemma~\ref{thm:area_law_uniform_separated_initial_marks} and
    (\ref{eq:area_law_repair_initial}), where $K_0=M\subseteq\mathbb R^2$ embeds by inclusion.
\end{proof}
% Source: 10-geometry.tex, sparse cells and initial marks, lines 212--233
% and 325--363; geometry:total-repairs, weighted sum, lines 683--692.
% Revision: adc7f1241b42e322a6451854ab7e4b4c146bf78a.
\begin{lemma}[Uniform weighted bound for the initial repair marks]
    \label{thm:area_law_initial_repair_weighted_sum}
    \lean{TNLean.PEPS.AreaLaw.Geometry.exists_uniform_initialRepairStage_weighted_sum_bound}
    \leanok
    \uses{def:area_law_repair_stage}
    For every $P\in\mathbb R$, there exists $B>0$ (depending only on $P$ and $\delta_0$) such that
    for any $\lambda\ge0$, $o\in\mathbb R^2$, finite $\Lambda\subseteq\mathbb Z^2$,
    $A\subseteq\Lambda$ with $b=|\partial_\Lambda A|$, $C_0\ge2$, $k_0\ge50{,}000{,}000$, and
    $w:K_0\to\mathbb R$ satisfying $w(p)\le\lambda(k+1)^P$ whenever $p\in K_0$, $k\ge k_0$, and
    $S_p^{(0)}\le2^k$,
    \begin{align}
        \sum_{p\in K_0}w(p)
         & \le576\lambda B(2C_0+1)^2b.
        \label{eq:area_law_initial_repair_weighted_sum}
    \end{align}
\end{lemma}

\begin{proof}
    \leanok
    \uses{thm:area_law_initial_repair_mark_facts,
        thm:area_law_belt_mark_count,thm:area_law_fine_belt_scales,
        thm:area_law_polynomial_finite_sums}
    Fix $B$ from Lemma~\ref{thm:area_law_polynomial_finite_sums} at exponent $P$. By
    Lemma~\ref{thm:area_law_initial_repair_mark_facts}, each $p\in K_0$ has a layer
    $\kappa(p)\ge k_0$ and $z\in F_{\kappa(p)}$ with $v_p^{(0)}\in M_{\ell_{\kappa(p)},z}$ and
    $S_p^{(0)}=2^{\ell_{\kappa(p)}}\le2^{\kappa(p)}$. Partitioning $K_0$ into
    $K_k=\{p\in K_0:\kappa(p)=k\}$ over the finite set $J$ of chosen layers,
    Lemma~\ref{thm:area_law_belt_mark_count} gives $|K_k|\le9|F_k|$ and hence
    \begin{align}
        \sum_{p\in K_0}w(p)
         & =\sum_{k\in J}\sum_{p\in K_k}w(p)
        \le\sum_{k\in J}9\lambda(k+1)^P|F_k|
        \le576\lambda(2C_0+1)^2b\sum_{k\in J}(k+1)^P2^{-\delta_0 k/2}
        \le576\lambda B(2C_0+1)^2b.\notag
    \end{align}
\end{proof}

% Source: 10-geometry.tex, geometry:belt-count, lines 220--233;
% initial fan runs, lines 299--318; geometry:total-repairs, lines 685--692.
% Revision: adc7f1241b42e322a6451854ab7e4b4c146bf78a.
\begin{lemma}[Uniform count of the initial belt regions]
    \label{thm:area_law_initial_belt_count}
    \lean{TNLean.PEPS.AreaLaw.Geometry.exists_uniform_initial_belt_card_bound}
    \leanok
    \uses{def:area_law_repair_stage,def:area_law_initial_region_indices,
        def:area_law_hamiltonian,def:area_law_belt_cells,
        def:area_law_fan_runs,def:area_law_belt_fan_color}
    There exists $B>0$ such that for any $o\in\mathbb R^2$, finite $\Lambda\subseteq\mathbb Z^2$,
    cut $A\subseteq\Lambda$, $C_0\ge2$, and $k_0\ge50{,}000{,}000$, the initial belt summand
    $\mathcal I_{\mathrm{belt}}=\coprod_{k\ge k_0}\coprod_{z\in F_k}\mathcal R_{k,z}$
    (Definition~\ref{def:area_law_initial_region_indices}) is finite and satisfies
    \begin{align}
        |\mathcal I_{\mathrm{belt}}|
         & \le B(2C_0+1)^2|\partial_\Lambda A|.
        \label{eq:area_law_initial_belt_count}
    \end{align}
\end{lemma}

\begin{proof}
    \leanok
    \uses{thm:area_law_initial_repair_mark_facts,def:area_law_fan_runs,
        thm:area_law_cell_fan_run_count,thm:area_law_polynomial_finite_sums}
    Take $B_0>0$ from Lemma~\ref{thm:area_law_polynomial_finite_sums} at $P=0$ and set $B=512B_0$.
    Since $J=\{k\ge k_0:F_k\ne\varnothing\}$ is finite
    (Lemma~\ref{thm:area_law_initial_repair_mark_facts}) and $|\mathcal R_{k,z}|\le8$
    (Lemma~\ref{thm:area_law_cell_fan_run_count}),
    \begin{align}
        |\mathcal I_{\mathrm{belt}}|
         & =\sum_{k\in J}\sum_{z\in F_k}|\mathcal R_{k,z}|
        \le8\sum_{k\in J}|F_k|
        \le512(2C_0+1)^2|\partial_\Lambda A|\sum_{k\in J}2^{-\delta_0 k/2}
        \le B(2C_0+1)^2|\partial_\Lambda A|.\notag
    \end{align}
\end{proof}
